\section{Invariant flags and the word algebra}

\begin{lemma}[Composition-series triangular form]
    \label{lem:asym_composition_gauge}
    \lean{exists_linearEquiv_blockTriangular_of_flag}
    \leanok
    Let $V$ be a finite-dimensional complex vector space, let
    $0=H_0\subseteq H_1\subseteq\cdots\subseteq H_r=V$ be a flag of
    subspaces, and let $(f_i)_i$ be endomorphisms of $V$ each preserving every
    $H_k$. Fix coordinates on each subquotient $H_k/H_{k-1}$. There is a
    linear isomorphism of $V$ with the graded coordinate space
    $\bigoplus_k\C^{\dim H_k/H_{k-1}}$ in which every $f_i$ is block upper
    triangular, its $k$-th diagonal block is the matrix of the map induced by
    $f_i$ on $H_k/H_{k-1}$, the coordinates of a vector of $H_k$ in block $k$
    are the coordinates of its class in $H_k/H_{k-1}$, and the coordinates of
    a vector of $H_{k-1}$ vanish in the blocks $k,k+1,\ldots,r$. In
    particular $\dim V=\sum_k\dim H_k/H_{k-1}$.
\end{lemma}

\begin{proof}
    \leanok
    Choose a complement of $H_{k-1}$ in $H_k$ for every $k$, and send a
    graded coordinate vector to the sum of the lifts of its blocks through
    these complements. If the image of a vector lies in $H_p$, its blocks past
    $p$ vanish: otherwise the last nonzero block, modulo the earlier ones,
    would give a nonzero class of a subquotient lying in $H_{p}$. Hence the
    map is injective, and comparison of dimensions makes it an isomorphism.
    Applying this vanishing to the images $f_i(H_k)\subseteq H_k$ gives block
    triangularity, and reading the $k$-th block modulo $H_{k-1}$ identifies
    the diagonal blocks.
\end{proof}

\begin{lemma}[Diagonal blocks of block-triangular products]
    \label{lem:asym_triangular}
    \lean{Matrix.blockDiag'_mul_of_blockTriangular, Matrix.blockTriangular_evalWord,
        Matrix.blockDiag'_evalWord, Matrix.blockProj_mul_blockEmbed_self,
        Matrix.blockProj_mul_blockEmbed_of_ne, Matrix.blockProj_mul_mul_blockEmbed}
    \leanok
    \uses{def:kraus_eval_word}
    Let $o$ be a finite index set with block sizes $n_k$, $k\in o$, and let
    $T^1,\ldots,T^d$ be block upper triangular for the resulting
    decomposition of $\MN{\sum_k n_k}$. Write $T^i_{kk}$ for the $k$-th
    diagonal block of $T^i$. Then $T^w$ is block upper triangular for every
    word $w$, and its $k$-th diagonal block is the word evaluation of the
    diagonal blocks,
    \begin{align}
        (T^w)_{kk} & =T^{i_1}_{kk}T^{i_2}_{kk}\cdots T^{i_N}_{kk}
        \label{eq:asym_triangular_word}
    \end{align}
    for $w=(i_1,\ldots,i_N)$. If $P_k$ is the coordinate embedding of the
    $k$-th block and $Q_k$ the coordinate projection onto it, then for every
    matrix $M$ on the same space, $Q_kMP_k=M_{kk}$, and
    \begin{align}
        Q_kP_l & =\begin{cases}\id, & k=l,\\0, & k\ne l.\end{cases}
        \label{eq:asym_block_biorthogonality}
    \end{align}
\end{lemma}

\begin{proof}
    \leanok
    For two block upper-triangular matrices $X,Y$,
    $(XY)_{kk}=\sum_l X_{kl}Y_{lk}$; the terms with $l<k$ vanish because
    $X_{kl}=0$, and those with $l>k$ vanish because $Y_{lk}=0$, leaving
    $(XY)_{kk}=X_{kk}Y_{kk}$. Inducting on the word length preserves block
    triangularity and extends this diagonal-block factorization to
    (\ref{eq:asym_triangular_word}). Since $P_k$ embeds the $k$-th coordinate
    block and $Q_k=P_k^{\mathsf T}$ projects onto it, direct coordinate
    multiplication gives $Q_kMP_k=M_{kk}$ and
    (\ref{eq:asym_block_biorthogonality}).
\end{proof}

\begin{lemma}[Products of strictly block-triangular matrices]
    \label{lem:wordflags_strict_triangular_nilpotent}
    \lean{Matrix.StrictBlockTriangular, Matrix.evalWord_eq_zero_of_strictBlockTriangular}
    \leanok
    \uses{def:kraus_eval_word}
    Let the blocks of a block decomposition carry labels in
    $\{0,\ldots,r-1\}$, and let $R^1,\ldots,R^d$ be strictly block upper
    triangular: every entry whose column label is at most its row label
    vanishes. Then $R^w=0$ for every word $w$ of length at least $r$.
\end{lemma}

\begin{proof}
    \leanok
    By induction on the length of $w$, the entry of $R^w$ in row block label
    $a$ and column block label $b$ vanishes unless $b\ge a+|w|$. For
    $|w|\ge r$ no label $b<r$ satisfies this.
\end{proof}

\begin{definition}[Word algebra and positive ideal]
    \label{def:asym_word_algebra}
    \lean{WordAlgebra, WordAlgebra.aug}
    \leanok
    Fix $d\geq1$. The \emph{word algebra} is the free unital associative
    algebra $F=\C\langle x_1,\ldots,x_d\rangle$, realized as the algebra of
    finite $\C$-linear combinations of words in the $d$ letters under
    concatenation, so that the words form a basis. The \emph{augmentation}
    $\mathrm{aug}:F\to\C$ is the algebra map sending every generator to
    zero; on a basis word $w$, $\mathrm{aug}(w)=1$ if $w$ is empty and $0$
    otherwise. The \emph{positive ideal} is the two-sided ideal
    $F_+=\ker(\mathrm{aug})$ of elements with zero constant term.
\end{definition}

\begin{definition}[Word traces]
    \label{def:asym_word_traces}
    \lean{WordAlgebra.traceChar, WordAlgebra.traceWord}
    \leanok
    \uses{def:asym_word_algebra}
    For a finite-dimensional $F$-module $M$, the \emph{trace character} is
    the linear functional $p\mapsto\tr(\rho_M(p))$ on $F$, and the
    \emph{word trace} of a word $w$ is $\tr(\rho_M(w))$.
\end{definition}

\begin{lemma}[Positive-ideal separation of linear functionals]
    \label{lem:asym_ker_aug_separation}
    \lean{WordAlgebra.linearMap_eq_on_ker_aug}
    \leanok
    \uses{def:asym_word_algebra}
    Two $\C$-linear functionals on $F$ agreeing on every nonempty word
    agree on the positive ideal $F_+$.
\end{lemma}

\begin{proof}
    \leanok
    Their difference agrees with $(\chi_1(1)-\chi_2(1))\cdot\mathrm{aug}$ on
    every basis word, hence on all of $F$ by linearity; on $F_+$ the
    augmentation vanishes, so the two functionals agree there.
\end{proof}

\begin{lemma}[The zero module is the only simple module with trivial
        generator action]
    \label{lem:asym_zero_module}
    \lean{WordAlgebra.finrank_eq_one_of_isSimpleModule_of_forall_smul_eq_zero}
    \leanok
    \uses{def:asym_word_algebra}
    Any finite-dimensional simple $F$-module on which every generator acts as
    zero has dimension one.
\end{lemma}

\begin{proof}
    \leanok
    If every generator acts as zero, $F$ acts through the augmentation, so
    the $F$-submodules and the $\C$-subspaces coincide; simplicity as a
    $\C$-vector space forces dimension one.
\end{proof}

\begin{lemma}[Nilpotent modules of the word algebra]
    \label{lem:wordflags_nilpotent_module}
    \lean{WordAlgebra.ofWord_smul_eq_zero_of_traceWord_eq_zero,
        WordAlgebra.traceWord_eq_zero_of_forall_ofWord_smul_eq_zero}
    \leanok
    \uses{def:asym_word_traces}
    Let $M$ be a finite-dimensional $F$-module. If the word trace of every
    nonempty word vanishes on $M$, then every word of length at least
    $\dim M$ acts on $M$ as zero. Conversely, if every word of some length
    acts as zero, the word trace of every nonempty word vanishes.
\end{lemma}

\begin{proof}
    \leanok
    The nonunital algebra generated by the letters is spanned by nonempty
    words, so every element has zero trace together with all its positive
    powers and is therefore nilpotent. A nonunital algebra of nilpotent
    endomorphisms of $M$ annihilates every product of $\dim M$ of its
    elements, and longer words contain such a product as a prefix.
    Conversely, a power of a nonempty word is a word of arbitrarily large
    length, so every nonempty word acts nilpotently and has trace zero.
\end{proof}

\begin{lemma}[Trace additivity over an invariant subspace]
    \label{lem:wordflags_trace_invariant_subspace}
    \lean{LinearMap.trace_eq_trace_restrict_add_trace_quotient}
    \leanok
    Let $f$ be an endomorphism of a finite-dimensional vector space $V$
    preserving a subspace $W$. Then the trace of $f$ is the trace of its
    restriction to $W$ plus the trace of the map induced on $V/W$.
\end{lemma}

\begin{proof}
    \leanok
    Choose a projection $P$ onto $W$ along a complement. Then
    $\tr f=\tr(fP)+\tr(f(\id-P))$; the first term is the trace of the
    restriction, and cyclicity of the trace along the quotient map and its
    section through the complement identifies the second with the trace on
    $V/W$.
\end{proof}

\begin{lemma}[A star-closed matrix algebra acts semisimply]
    \label{lem:wordflags_star_closed_semisimple}
    \lean{StarSubalgebra.isSemisimpleModule}
    \leanok
    A unital algebra of complex square matrices closed under conjugate
    transposition makes the coordinate space a semisimple module: every
    invariant subspace has an invariant complement.
\end{lemma}

\begin{proof}
    \leanok
    The orthogonal complement of an invariant subspace is invariant, because
    the conjugate transposes of the members of the algebra belong to it and
    preserve the subspace.
\end{proof}

\begin{lemma}[Factoring a linear map through a surjection on a subspace]
    \label{lem:wordflags_subspace_factorization}
    \lean{Matrix.exists_mulVec_eq_of_surjOn}
    \leanok
    Let $U$ be a subspace of a complex vector space, let $f$ map $U$ onto
    $\C^m$, and let $g$ be a linear map to $\C^n$ vanishing on the vectors of
    $U$ that $f$ sends to zero. Then there is an $n\times m$ matrix $X$ with
    $g(x)=Xf(x)$ for every $x\in U$.
\end{lemma}

\begin{proof}
    \leanok
    Take $X$ to be $g$ composed with a linear right inverse of $f$ on $U$;
    the difference of $x$ and the lift of $f(x)$ lies in the kernel of $f$,
    so $g$ takes the same value on both.
\end{proof}

\begin{lemma}[One step against another]
    \label{lem:asym_uniq_step}
    \lean{CoveringMatching.step_sup_eq_of_lt,
        CoveringMatching.step_sup_eq_of_lt',
        CoveringMatching.step_inf_eq_of_lt,
        CoveringMatching.step_lt_symm}
    \leanok
    In a modular lattice let $P\lessdot Q$ and $P'\lessdot Q'$ be covering
    pairs and put $U=Q\wedge Q'$. If $P\vee(Q\wedge P')<P\vee(Q\wedge Q')$,
    then $P\vee U=Q$, $Q\wedge P'\leq P$, $P'\vee U=Q'$,
    $Q'\wedge P\leq P'$, $P\wedge Q'=Q\wedge P'$, and the strict inequality
    holds with the two pairs exchanged.
\end{lemma}

\begin{proof}
    \leanok
    $P\vee(Q\wedge P')$ and $P\vee U$ lie between $P$ and $Q$ with the first
    strictly below the second, so the first is $P$ and the second is $Q$,
    since nothing lies strictly between $P$ and $Q$. Then $P'\vee U$ equals
    $P'$ or $Q'$, and the first case gives $U\leq Q\wedge P'\leq P$, so
    $P\vee U=P\neq Q$; hence $P'\vee U=Q'$. Next $P'\vee(Q'\wedge P)$ equals
    $P'$ or $Q'$; in the second case $Q'\leq P\vee P'$ and the modular law
    gives $U\leq Q\wedge(P\vee P')=P\vee(Q\wedge P')=P$, again contradicting
    $P\vee U=Q$. The two inclusions $Q\wedge P'\leq P\wedge Q'$ and
    $Q'\wedge P\leq Q\wedge P'$ give the equality, and the exchanged
    inequality is $P'<Q'$.
\end{proof}

\begin{proposition}[The matching]
    \label{prop:asym_uniq_matching}
    \lean{CoveringMatching.CoveringChain,
        CoveringMatching.CoveringChain.jump,
        CoveringMatching.CoveringChain.jump_spec,
        CoveringMatching.CoveringChain.jump_jump,
        CoveringMatching.CoveringChain.jumpEquiv,
        CoveringMatching.CoveringChain.length_eq,
        CoveringMatching.CoveringChain.jump_sup_eq}
    \leanok
    \uses{lem:asym_uniq_step}
    Let $\bot=M_0\lessdot M_1\lessdot\cdots\lessdot M_r=\top$ and
    $\bot=M'_0\lessdot\cdots\lessdot M'_{r'}=\top$ be two chains of covering
    pairs in a modular lattice. For every $k$ there is exactly one $k'$ with
    $M_{k-1}\vee(M_k\wedge M'_{k'-1})<M_{k-1}\vee(M_k\wedge M'_{k'})$, namely
    the least $k'$ with $M_k\wedge M'_{k'}\not\leq M_{k-1}$. Writing
    $k'=\pi(k)$, the map $\pi$ is a bijection with inverse the same rule for
    the chains exchanged, so $r=r'$; and with $U=M_k\wedge M'_{\pi(k)}$ and
    $K=M_{k-1}\wedge M'_{\pi(k)}$ one has $M_{k-1}\vee U=M_k$,
    $M'_{\pi(k)-1}\vee U=M'_{\pi(k)}$, $U\wedge M_{k-1}=U\wedge
        M'_{\pi(k)-1}=K$ and $K=M_k\wedge M'_{\pi(k)-1}$.
\end{proposition}

\begin{proof}
    \leanok
    The elements $M_{k-1}\vee(M_k\wedge M'_j)$ increase with $j$ from
    $M_{k-1}$ to $M_k$, so exactly one step is strict, and the least $k'$
    above is such a step. Lemma~\ref{lem:asym_uniq_step} applied to the two
    covering pairs shows that any strict step equals it, that the rule for
    the exchanged chains sends $k'$ back to $k$, and gives the identities.
\end{proof}
